\input format.tex

%\usepackage{arydshln}

\begin{document}

%%%
\setlength{\arrayrulewidth}{0.5pt}
\fontsize{9.3pt}{17pt}\selectfont
\color{gray2}

\vspace*{0mm}

\bigskip
\noindent
基于本次您的检测结果，建议您从以下几个方面进行干预：
\begin{compactenum}[\mbox{1、}]
\xiaowuhao \color{gray2} \lantxh {\item[1、]
鉴于您的肠道菌群多样性低，建议多补充益生元和益生菌，帮助调节肠道菌群平衡；您体内较高含量的致病菌可能诱发相关肠道疾病，建议您持续监测肠道健康，并到正规医院咨询专科医生，遵从医生建议做进一步检查。
}
\xiaowuhao \color{gray2} \lantxh {\item[2、]
%您的肠道菌群可能增加您患2型糖尿病、结直肠癌、炎症性肠病的风险，建议您-饮食，少吃-的食物，补充益生菌，定期监测肠道健康，积极预防。
您的肠道菌群可能增加您患2型糖尿病、结直肠癌、炎症性肠病的风险，建议您控制体重、保持大便通畅、注意饮食卫生、避免摄入不耐受的食物及损伤肠粘膜的药物，定期监测肠道健康，积极预防。
}
\xiaowuhao \color{gray2} \lantxh {\item[3、]
您肠道内参与脂肪总量、甘油三酯代谢、鞘脂代谢等的菌群指标异常，
可能导致动脉硬化、心脑血管疾病、便秘等疾病，建议您关注血糖血脂血压、保持心态平和、养成定时排便的习惯、适量多饮水、戒烟酒、避免熬夜，定期监测肠道健康。
}
\end{compactenum}

\bigskip
\tabcolsep=6pt
\fontsize{9.3pt}{9.8pt}\selectfont
\begin{tctabularx}{tabularx={p{15.4cm}<{\centering}}}
\\[-6.769pt]
\cellcolor{kelly} \raisebox{1pt}{\color{black70}\xiaosihao\lantxh\bf 膳食方案} \\
\end{tctabularx}

{
\lantxh
\vspace*{-4.25mm}
\noindent
\begin{longtable}{|p{1.25cm}@{}p{2cm}|p{11.732cm}|}
\hline
\parbox[c]{\hsize}{\vskip7pt\raisebox{-1.5ex}{\includegraphics[width=.9cm]{zhushi.pdf}}\vskip7pt} & \parbox[c]{\hsize}{\centering\xiaowuhao \color{gray2}\vskip7pt主食\vskip7pt} & \parbox[c]{\hsize}{\xiaowuhao \color{gray2}\vskip7pt推荐高膳食纤维饮食。可多食用富含膳食纤维的谷物，如糙米饭、燕麦饭、全麦面、玉米馒头等。\vskip7pt}\\
\hline
\parbox[c]{\hsize}{\vskip7pt\raisebox{-1.5ex}{\includegraphics[width=.9cm]{roulei.pdf}}\vskip7pt} & \parbox[c]{\hsize}{\centering\xiaowuhao \color{gray2}\vskip7pt肉类\vskip7pt} & \parbox[c]{\hsize}{\xiaowuhao \color{gray2}\vskip7pt选择吃瘦肉，避免高脂食物，如肥猪肉、肥牛、猪脑等，适量减少摄入五花肉、排骨等。\vskip7pt}\\
\hline
\parbox[c]{\hsize}{\vskip7pt\raisebox{-1.5ex}{\includegraphics[width=.9cm]{yuxialei.pdf}}\vskip7pt} & \parbox[c]{\hsize}{\centering\xiaowuhao \color{gray2}\vskip7pt鱼虾类\vskip7pt} & \parbox[c]{\hsize}{\xiaowuhao \color{gray2}\vskip7pt宜多吃鱼虾类优质蛋白，宜采用清蒸烹饪方法。\vskip7pt}\\
\hline
\parbox[c]{\hsize}{\vskip7pt\raisebox{-1.5ex}{\includegraphics[width=.9cm]{danlei.pdf}}\vskip7pt} & \parbox[c]{\hsize}{\centering\xiaowuhao \color{gray2}\vskip7pt蛋奶、大豆类\vskip7pt} & \parbox[c]{\hsize}{\xiaowuhao \color{gray2}\vskip7pt每天一个鸡蛋，奶类宜选用酸奶，每天适量多摄入豆类及豆制品，如黄豆、豆腐、豆浆等。\vskip7pt}\\
\hline
\parbox[c]{\hsize}{\vskip7pt\raisebox{-1.5ex}{\includegraphics[width=.9cm]{shucaishuiguo.pdf}}\vskip7pt} & \parbox[c]{\hsize}{\centering\xiaowuhao \color{gray2}\vskip7pt蔬菜水果\vskip7pt} & \parbox[c]{\hsize}{\xiaowuhao \color{gray2}\vskip7pt多吃富含膳食纤维的新鲜果蔬，如柚子、柑橘、苹果、香蕉、火龙果、洋葱、芹菜、洋白菜、豌豆、蚕豆等。\vskip7pt}\\
\hline
\parbox[c]{\hsize}{\vskip7pt\raisebox{-1.5ex}{\includegraphics[width=.9cm]{you.pdf}}\vskip7pt} & \parbox[c]{\hsize}{\centering\xiaowuhao \color{gray2}\vskip7pt食用油、盐\vskip7pt} & \parbox[c]{\hsize}{\xiaowuhao \color{gray2}\vskip7pt以橄榄油、菜籽油作为膳食油脂的主要来源，食用油每天不超过25g，食盐摄入量每天不超过6g。\vskip7pt}\\
\hline
\parbox[c]{\hsize}{\vskip7pt\raisebox{-1.5ex}{\includegraphics[width=.9cm]{other.pdf}}\vskip7pt} & \parbox[c]{\hsize}{\centering\xiaowuhao \color{gray2}\vskip7pt其他\vskip7pt} & \parbox[c]{\hsize}{\xiaowuhao \color{gray2}\vskip7pt蔬菜的烹饪方式以轻烫为宜，不宜过度加工。如患痛风、溃疡性结肠炎、慢性腹泻、食物不耐受等有饮食禁忌的疾病，请优先遵循疾病的饮食原则。\vskip7pt}\\
\hline
\end{longtable}
\vspace*{-6mm}
\begin{longtable}{p{0.1cm}p{15.1cm}}
\includegraphics{xiaofangkuai.pdf} & {\qihao 以上膳食方案仅供参考，您的实际饮食请结合自身具体情况调整。} \\
\end{longtable}

}
\end{document}
